
\documentclass[11pt, a4paper]{article}

% ===== ESSENTIAL PACKAGES =====
\usepackage{geometry}
\geometry{a4paper, margin=1.2in}
\usepackage{setspace}
\usepackage{array}
\usepackage{xcolor}
\onehalfspacing
\setlength{\parskip}{1em}
\usepackage{multicol}

\usepackage{catchfile}

% ===== WORD COUNT COMMAND =====
\newcommand{\wordcount}{%
    \ifnum\pdf@shellescape=1%
        \immediate\write18{texcount -1 -sum Oblig1.tex > wordcount.txt}%
        \CatchFileDef{\words}{wordcount.txt}{}%
        \words%
    \else%
        [Enable -shell-escape]%
    \fi%
}
\usepackage{mdframed}

\newmdenv[
    leftmargin=-50pt,
    rightmargin=0pt,
    innerleftmargin=3pt,
    innerrightmargin=0pt,
    innertopmargin=5pt,
    innerbottommargin=5pt,
    skipabove=5pt,
    skipbelow=5pt,
    linewidth=1pt,
    linecolor=black,
    topline=true,
    rightline=false,
    bottomline=false
]{sectionboxinner}

\newenvironment{sectionbox}[1][]{%
    \begin{sectionboxinner}
    \textbf{\large #1}
    \vspace{1pt}
    \newline
}{%
    \end{sectionboxinner}
}

\newmdenv[
    leftmargin=0pt,
    rightmargin=0pt,
    innerleftmargin=0pt,
    innerrightmargin=0pt,
    innertopmargin=3pt,
    innerbottommargin=3pt,
    skipabove=8pt,
    skipbelow=8pt,
    linewidth=1.2pt,
    linecolor=gray!60,  % Lighter gray color
    topline=false,
    rightline=false,
    bottomline=false
]{subsectionboxinner}

\newenvironment{subsectionbox}[1][]{%
    \begin{subsectionboxinner}
    \textbf{#1}
    \vspace{3pt}
    \newline
}{%
    \end{subsectionboxinner}
}

% ===== CUSTOM ENVIRONMENTS FOR PHILOSOPHICAL ARGUMENTS =====
\newenvironment{argument}{\begin{quote}\itshape\textbf{Argument: }}{\end{quote}}
\newenvironment{objection}{\begin{quote}\itshape\textbf{Innvending: }}{\end{quote}}
\newenvironment{reply}{\begin{quote}\itshape\textbf{Reply: }}{\end{quote}}

% ===== DOCUMENT BEGIN =====
\title{Kap. 17 - Den sosiale kontrakten - hvorfor trenger vi stat og myndighet?}
\author{Oliver Ekeberg}
\date{\today}

\begin{document}
\maketitle

\tableofcontents


\begin{sectionbox}[Antroposentrisk miljøetikk]

Ta en fjellreven. Et vakkert vesen som er søt, og har en rolle i økosystemet. Har det noe å si for oss at den er utryddet? Er det moralsk at den blir utryddet på grunn av mennesker, enten ved jakt eller menneskeskapte klimaforandringer? En filosofisk posisjon er at naturen ligger utenfor moralens område. Moral dreier seg om mennesker og menneskelig handling, og det er personer som er målet på moralsk verdi. Dermed må et moralsk argument for dyr eller miljø være Antroposentrisk

\begin{quote}
    \textbf{\textit{Antroposentrisk:}} å ødelegge eller skade dyr eller miljø har moralsk verdi på grunn av verdien naturen har for mennesket.
\end{quote}

En person som spiser kjøtt produsert av industrielt landbruk som påfører dyr lidelse, er ikke i seg selv slem, men han er ikke barmhjertig. Siden dyden barmhjertighet utelukker at man er likegyldig overfor andre veseners lidelser.\\

Dette er et antroposentrisk argument fordi det tillegger naturen \textit{instrumental verdi}. Vi har moralske grunner som tar hensyn til dyr eller natur fordi det er godt for oss. Det er galt å skape ubalanse i økosystemet fordi det er skadelig for oss. Mange vil påstå at en slik holdning instrumentell holdning er årsaken til problemene våre. Trenger vi en miljøetikk basert på at \textit{naturen har iboende verdi?}

\end{sectionbox}

\end{document}